\section{Conclusiones}

Este artículo ha analizado el marco teórico y propuestas para el Portal Agro Comercial del Huila, demostrando cómo las tecnologías digitales pueden reducir significativamente la intermediación en la comercialización agrícola. Los 19 artículos base, junto con literatura adicional reciente, proporcionan una base sólida para implementar un portal que empodere a los productores y conecte directamente con mercados.

La integración de IoT, e-commerce y plataformas tecnológicas aborda problemas estructurales del Huila, como brechas digitales y dependencia de intermediarios. Las propuestas de aplicación ofrecen soluciones prácticas, con mejoras innovadoras que incluyen IA, blockchain y capacitación inclusiva.

Para el departamento del Huila, el portal representa una oportunidad para aumentar la competitividad agrícola, generar empleos en tecnología y promover sostenibilidad. Se recomienda políticas públicas de apoyo, inversiones en infraestructura digital y alianzas con instituciones académicas para escalabilidad.

El análisis comparativo revela consensos en la necesidad de adopción gradual y contextualizada, diferenciando enfoques técnicos de comportamentales. Futuras investigaciones podrían evaluar el impacto cuantitativo del portal en indicadores económicos del Huila.

En conclusión, el Portal Agro Comercial no solo reduce intermediación, sino que transforma la agricultura del Huila hacia un modelo digital, inclusivo y sostenible, contribuyendo al desarrollo regional y nacional.